Elementos que debe contener una etiqueta:

\begin{itemize}
    \item Nombre de la marca (y eslogan, de aplicar)
    \item Sobrenombre del producto
    \item Nombre del producto (Descripción)
    \item Sellos (NOM-)
        \begin{itemize}
	        \item Sodio
	        \item Calorias
	        \item Grasas
	        \item Azúcares añadidos (no naturales)
        \end{itemize}
    
    \item Ingredientes (De mayor a menor)
    \item Información nutrimental (por porciones)
\end{itemize}

\definition{Trazabilidad}{Información para identificar el proceso o lote donde se generó el producto y todos los ingredientes involucrados.}

Diferencia entre fecha de caducidad y fecha de consumo preferente:

\begin{itemize}
    \item Fecha de caducidad: Fecha en que el producto comienza a poder causar daño al consumirse.
    \item Fecha de consumo preferente: Fecha en que el producto pierde sus propiedades organolépticas, pero aún así es seguro de consumir.
\end{itemize}

\section{Dependencias de seguridad alimentaria}

\subsection{Nivel nacional}

Las dependencias de gobierno responsables de detectar fraudes alimentarios son:

\begin{itemize}
    \item COFEPRIS y COEPRIS
    \item PROFECO
    \item SSA\footnote{Esta opera solo en caso de presentarse casos de enfermedades transmitidas por alimentos (ETA)}
\end{itemize}

\subsection{Nivel internacional}

Dependencias del Codex Alimentarius

\begin{itemize}
    \item AOAC: Dependencia encargada de generar métodos oficiales para alimentos, bebidas, cosméticos, etcétera.
    \item ACC: Específico de cereales
\end{itemize}

Concepto de puntos frios y centros de calor:

\begin{itemize}
    \item Punto frio: Lugar por donde se escapa más fácilmente el calor
    \item Centro de calor: Lugar que retiene mejor el calor
\end{itemize}

Dato: La proteína suele ser el ingrediente más caro en la producción de un alimento.
